\subsection{Clasificaci\'on supervisada}
Se entrenaron dos modelos supervisados sobre los 64 coeficientes MFCC estandarizados. El primer modelo correspondi\'o a una red MLP implementada en TensorFlow/Keras con topolog\'ia $64 \rightarrow 128 \rightarrow 64 \rightarrow 5$. En las capas ocultas se aplicaron Batch Normalization, activaci\'on ReLU y Dropout con tasas 0.30 y 0.20, respectivamente. La funci\'on de p\'erdida utilizada fue entrop\'ia cruzada categ\'orica dispersa, optimizada mediante Adam.

El segundo modelo correspondi\'o a un ensamble por votaci\'on suave compuesto por Random Forest, Extra Trees e HistGradientBoosting. Se utilizaron probabilidades posteriores estimadas por cada componente y se agregaron mediante votaci\'on probabil\'istica. La red neuronal fue entrenada durante 69 \'epocas efectivas con parada temprana.

\begin{table}[H]
\centering
\caption{Desempe\~no de componentes del ensamble en validaci\'on.}
\begin{tabular}{lrr}
\hline
Modelo & F1 ponderado validaci\'on & Tiempo (s) \\
\hline
random_forest & 0.443 & 2.471 \\
extra_trees & 0.422 & 1.326 \\
hist_gradient_boosting & 0.422 & 13.987 \\
soft_voting_ensemble & 0.437 & 16.853 \\
\hline
\end{tabular}
\end{table}

\begin{table}[H]
\centering
\caption{Comparaci\'on de clasificadores sobre el conjunto de prueba.}
\begin{tabular}{lrrr}
\hline
Modelo & Accuracy & F1 macro & F1 ponderado \\
\hline
MLP TensorFlow & 0.451 & 0.451 & 0.450 \\
Soft Voting Ensemble & 0.491 & 0.428 & 0.469 \\
\hline
\end{tabular}
\end{table}

El F1-score fue priorizado porque el problema contiene cinco especies y la distribuci\'on de clases puede inducir diferencias entre desempe\~no global y desempe\~no por clase. Las matrices de confusi\'on permitieron inspeccionar errores de asignaci\'on entre especies ac\'usticamente cercanas, mientras que las curvas de aprendizaje permitieron controlar el sobreajuste asociado a la red neuronal.
